\subsection{THE RATIONAL LINE}

Consider the set \(\mathfrak{F}\) of symmetric open intervals \(]-a, +a[\) , where \(a\) runs through the set of rational numbers \(>0\) ; we shall show that \(\mathfrak{F}\) is a fundamental system of neighbourhoods of \(o\) in a topology compatible with the additive group structure of \(\mathbf{Q}\) .

The group Q is commutative, and axiom \((\mathbf{GV}_{\mathrm{II}}^{\prime})\) is clearly satisfied; it is therefore enough to show that axiom \((\mathbf{GV}_{\mathrm{I}}^{\prime})\) is also satisfied, in other words, that for each a \textgreater{} 0 there exists b \textgreater{} 0 such that the conditions \(|x| < b\) and \(|y| < b\) together imply \(|x + y| < a\) . The triangle inequality shows that we may take b = a/2.

DEFINITION 1. The rational line is the topological space consisting of the set Q together with the additive group topology for which the symmetric open intervals \(]-a, +a[\) (a \textgreater{} o) form a fundamental system of neighbourhoods of o.

The topological group Q thus defined is called the additive group of the rational line.

If \(a\) is any rational number \(> 0\) , there is an integer \(n > 0\) such that \(1 / n < a\) ; hence the open intervals \(] - \frac{1}{n}, + \frac{1}{n} [ (n = 1, 2, \ldots) \text{ form}\)

a fundamental system of neighbourhoods of o on the rational line.

We obtain a fundamental system of neighbourhoods of any point \(x \in Q\) by taking the open intervals \(]x - a, x + a[\) , where a runs through the set of rational numbers \textgreater0 (or the set of numbers \(1/n\) ).

Definition 1 is therefore equivalent to that given in Chapter I, § 1, no. 2.

For each pair of rational numbers \((a, b)\) such that a \textless{} b, there exists \(c \in Q\) such that a \textless{} c \textless{} b {[}for example \(c = (a + b)/2\) {]}; it follows that the rational line is a non-discrete Hausdorff space.

For each \(a > 0\) , let \(\mathbf{U}_a\) be the set of pairs \((x, y)\) in \(\mathbf{Q} \times \mathbf{Q}\) such that \(|x - y| < a\) . As \(a\) runs through the set of rational numbers \(> 0\) (or just the set of numbers \(1/n\) ), the sets \(\mathbf{U}_a\) form a fundamental system of entourages of the uniformity of the additive group \(\mathbf{Q}\) of the rational line. Relations (2) and (3) show that \(|x|, x^+\) and \(x^-\) are uniformly continuous on \(\mathbf{Q}\) . It follows that the functions \(\sup (x, y)\) and \(\inf (x, y)\) are uniformly continuous on \(\mathbf{Q} \times \mathbf{Q}\) .

